\documentclass[12pt]{article}
\usepackage[margin=1in]{geometry}
\usepackage{booktabs}
\usepackage{longtable}
\usepackage{caption}
\DeclareCaptionLabelSeparator{emdash}{---}
\captionsetup{font=normalsize,labelfont=sc,textfont=sc,labelsep=emdash}
\begin{document}
\begingroup
\footnotesize
\renewcommand{\arraystretch}{0.86}
\setlength{\tabcolsep}{4pt}
\setlength{\LTcapwidth}{\textwidth}
\begin{longtable}{l*{6}{c}}
\caption{Summary Statistics, Discrete-Dose Samples}\\
\toprule
 & \multicolumn{6}{c}{Panel A: WAZ, age \ensuremath{\le}10 in 2005 and 2009}\\
\cmidrule(lr){2-7}
                    &        Mean&   Std. dev.&        Min.&      Median&        Max.&        Obs.\\
\midrule
\textit{Outcomes (MxFLS-2, 2005)}&            &            &            &            &            &            \\
Weight-for-age z-score (WAZ)&       -0.30&        1.42&       -5.07&       -0.36&        4.51&       1,045\\
\addlinespace[6pt]\textit{Outcome change, 2005 to 2009}&            &            &            &            &            &            \\
Change in WAZ       &        0.55&        1.41&       -4.88&        0.46&        9.10&       1,045\\
\addlinespace[6pt]\textit{Dose levels (CHIRPS)}&            &            &            &            &            &            \\
2009 rainfall, z    &       -0.50&        0.83&       -2.24&       -0.39&        0.96&       1,045\\
Near normal (\ensuremath{-1<z<1})&        0.65&        0.48&        0.00&        1.00&        1.00&       1,045\\
Moderately dry (\ensuremath{-1.5<z\le -1})&        0.21&        0.41&        0.00&        0.00&        1.00&       1,045\\
Severely dry or worse (\ensuremath{z\le -1.5})&        0.14&        0.35&        0.00&        0.00&        1.00&       1,045\\

\bottomrule
\addlinespace[12pt]
\toprule
 & \multicolumn{6}{c}{Panel B: HAZ and BMIZ, age \ensuremath{\le}15 in 2005 and 2009}\\
\cmidrule(lr){2-7}
                    &        Mean&   Std. dev.&        Min.&      Median&        Max.&        Obs.\\
\midrule
\textit{Outcomes (MxFLS-2, 2005)}&            &            &            &            &            &            \\
Height-for-age z-score (HAZ)&       -0.58&        1.49&       -5.46&       -0.63&        5.65&       2,351\\
BMI-for-age z-score (BMIZ)&        0.36&        1.43&       -4.95&        0.30&        4.86&       2,310\\
\addlinespace[6pt]\textit{Outcome change, 2005 to 2009}&            &            &            &            &            &            \\
Change in HAZ       &        0.07&        1.37&       -5.93&        0.07&        6.35&       2,351\\
Change in BMIZ      &        0.23&        1.31&       -5.90&        0.17&        7.59&       2,295\\
\addlinespace[6pt]\textit{Dose levels (CHIRPS)}&            &            &            &            &            &            \\
2009 rainfall, z    &       -0.55&        0.84&       -2.24&       -0.53&        0.96&       2,351\\
Near normal (\ensuremath{-1<z<1})&        0.64&        0.48&        0.00&        1.00&        1.00&       2,351\\
Moderately dry (\ensuremath{-1.5<z\le -1})&        0.21&        0.41&        0.00&        0.00&        1.00&       2,351\\
Severely dry or worse (\ensuremath{z\le -1.5})&        0.16&        0.36&        0.00&        0.00&        1.00&       2,351\\

\bottomrule
\addlinespace[6pt]\multicolumn{7}{p{0.97\textwidth}}{\scriptsize \textit{Notes:} This table is for the discrete-dose specification, the main model. Obs. counts children. The sample is children with the outcome measured in MxFLS-2 (2005) and MxFLS-3 (2009), under the panel's age cutoff
in both survey waves, with a 2005 municipality. They were not measured in MxFLS-3 before November 2009. I use November 2009 as the cutoff because the 2009 rainfall deficit had not fully occurred before then. Dropping them from both groups equalizes treated and control measurement timing. Children in wet municipalities (z of 1 or above) are excluded.
Panel A uses the WAZ sample, Panel B the HAZ and BMIZ sample. The BMIZ sample is marginally smaller than the HAZ sample shown in Panel B, at 2,300 children. Statistics are at the 2005 baseline, except the change in WAZ, HAZ and BMIZ from 2005 to 2009, the outcome of the discrete-dose regressions, and 2009 rainfall. z is 2009 May--October CHIRPS rainfall, the maize growing season (Jessoe, Manning, and Taylor 2018), in the child's 2005 municipality in standard deviations
from its 1981--2004 mean, grouped into dose levels with the SPI cutoffs (McKee, Doesken, and Kleist 1993; WMO 2012). The 2009 rainfall z is continuous. The three dose-group indicators have means equal to the share of children in each dose group. The discrete-dose regressions have no controls, as the contdid package does not yet accommodate covariates (Callaway, Goodman-Bacon, and Sant'Anna 2026).}\\
\end{longtable}
\endgroup
\end{document}
